\section*{الفصل \ref{ch:g9:climate-past-and-present} --- مناخ الأرض بين الماضي والحاضر}

\begin{solution}{exo:g9:climate-past-and-present:1}
(أ) نحو \qty{-19}{\degreeCelsius}. (ب) نحو \qty{15}{\degreeCelsius}.
(ج) $15 - (-19) = 34$: أي نحو \qty{34}{\degreeCelsius}. (د) اثنان مما يلي:
ثاني أكسيد الكربون، وبخار الماء، والميثان.
\end{solution}

\begin{solution}{exo:g9:climate-past-and-present:2}
(أ) \omterm{def:g1:watching-the-weather:weather}{طقس}. (ب) \omterm{def:g5:weather-and-climate:climate}{مناخ}. (ج) \omterm{def:g5:weather-and-climate:climate}{مناخ} (وتغير \omterm{def:g5:weather-and-climate:climate}{مناخ}: تغيّر على مدى
عقود). (د) \omterm{def:g1:watching-the-weather:weather}{طقس}.
\end{solution}

\begin{solution}{exo:g9:climate-past-and-present:3}
(أ) \qty{19}{cm}، وإحدى عشرة سنة. (ب) $19 \div 11 \approx \qty{1.7}{cm}$.
(ج) يسقط \omterm{def:g4:clouds-rain-snow:precipitation}{ثلج} كل سنة فوق \omterm{def:g4:clouds-rain-snow:precipitation}{ثلج} السنوات السابقة، فكلما كانت الطبقة أعمق، كان
سقوطها أبكر.
\end{solution}

\begin{solution}{exo:g9:climate-past-and-present:4}
(أ) $22.4 - 21.0 = +\qty{1.4}{\degreeCelsius}$. (ب) $20.3 - 21.0 =
\qty{-0.7}{\degreeCelsius}$. (ج) \qty{0}{\degreeCelsius}. (د) أن ذلك الشهر
كان أبرد من متوسطه على مدى 1991--2020.
\end{solution}

\begin{solution}{exo:g9:climate-past-and-present:5}
(أ) من عقد 1880--1889 إلى عقد 1930--1939، ثم عقدا 1950--1959 و1960--1969.
(ب) نحو \qty{0.25}{\degreeCelsius} و\qty{0.80}{\degreeCelsius}. (ج) $0.80 -
0.25 = \qty{0.55}{\degreeCelsius}$. (د) لا: فالعقود من 1880--1889 إلى
1910--1919 صارت أبرد قليلًا، وكان عقد 1950--1959 أبرد قليلًا من عقد
1940--1949؛ لكن منذ عقد 1970--1979 صار كل عقد أدفأ من الذي
قبله.
\end{solution}

\begin{solution}{exo:g9:climate-past-and-present:6}
(أ) $(324.62 - 315.98) \div 10 \approx \qty{0.86}{ppm}$ في السنة. (ب)
$(427.35 - 401.01) \div 10 \approx \qty{2.63}{ppm}$ في السنة. (ج) $2.63
\div 0.86 \approx 3.1$: أي أسرع بنحو ثلاث مرات. (د) $(427.35 - 315.98)
\div 315.98 \times 100 \approx 35~\%$.
\end{solution}

\begin{solution}{exo:g9:climate-past-and-present:7}
(أ) $8 \times 10^5$ سنة. (ب) قبل نحو \num{120000} إلى \num{130000} سنة:
فمنحنى درجة الحرارة يرتفع فوق 0 هناك. (ج) نحو \qty{275}{ppm}. (د) نعم: فقمم
المنحنيين وقيعانهما تأتي في الأوقات نفسها.
\end{solution}

\begin{solution}{exo:g9:climate-past-and-present:8}
(أ) $3.2 \times 10 = \qty{32}{mm} = \qty{3.2}{cm}$. (ب) $2100 - 2025 =
75$ سنة؛ $75 \times 3.2 = \qty{240}{mm}$، أي نحو \qty{24}{cm} (وأكثر إذا
تسارع الارتفاع). (ج) ماء \omterm{def:g2:water-on-earth:ocean}{البحر} يتمدد حين يسخن؛ والجليد على اليابسة يذوب فيسيل
ماؤه إلى \omterm{def:g2:water-on-earth:ocean}{البحر}. (د) لا: \omterm{def:g8:ice-on-earth:sea-ice}{فالجليد البحري} يطفو على \omterm{def:g2:water-on-earth:ocean}{البحر} أصلًا، فذوبانه لا
يرفعه.
\end{solution}

\begin{solution}{exo:g9:climate-past-and-present:9}
(أ) 90. (ب) 66. (ج) 99. (د) \emph{مرجح} تعني 66 فرصة من 100 على الأقل، لا
اليقين: فالصحيفة تجعل العبارة تبدو أوثق مما قاله
العلماء.
\end{solution}

\begin{solution}{exo:g9:climate-past-and-present:10}
الشتاء البارد الواحد في مكان واحد \emph{\omterm{def:g1:watching-the-weather:weather}{طقس}}: فهو يخبر بما حدث خلال بضعة
أشهر في بقعة واحدة. أما \omterm{def:g9:climate-past-and-present:global-warming}{الاحترار العالمي} فتغيّر في \emph{\omterm{def:g5:weather-and-climate:climate}{المناخ}}، يُقاس على
مدى عقود، وهو \emph{عالمي}: أي إنه المتوسط على الأرض كلها. وقد يحدث شتاء بارد
هنا بينما الأرض في مجملها أدفأ من قبل، ومتوسطات العقود تواصل
الارتفاع.
\end{solution}

\begin{solution}{exo:g9:climate-past-and-present:11}
(أ) $280 - 190 = \qty{90}{ppm}$، في نحو $2 \times 10^4$ سنة. (ب) $427
- 280 = \qty{147}{ppm}$، في $2025 - 1750 = 275$ سنة. (ج) الارتفاع الثاني
أكبر (147 مقابل 90~ppm). المعدلان: $90 \div (2 \times 10^4) =
4.5 \times 10^{-3}$~ppm في السنة، مقابل $147 \div 275 \approx 5.3 \times
10^{-1}$~ppm في السنة: أي أسرع بنحو 120 مرة. (د) أنه كان في 2019 أعلى منه
في أي وقت خلال مليوني سنة على الأقل (بثقة عالية).
\end{solution}

\begin{solution}{exo:g9:climate-past-and-present:12}
(أ) متساويان تقريبًا: 1.07 مقابل \qty{1.06}{\degreeCelsius}. فأفضل تقدير يقول
إن التأثير البشري يفسر نحو الاحترار المرصود كله. (ب) لا: فالأسباب الطبيعية
مجتمعة لم تغير درجة الحرارة إلا بمقدار \qty{0.1}{\degreeCelsius} على الأكثر، أي
أقل بعشر مرات من الاحترار المرصود. (ج) 66 فرصة من 100 على الأقل لأن يقع
الاحترار الناجم عن الإنسان بين 0.8 و\qty{1.3}{\degreeCelsius}.
\end{solution}

\begin{solution}{pb:g9:climate-past-and-present:1}
\textbf{1.} لأنها الفترة المرجعية: فالشذوذات تُقاس انطلاقًا من متوسطها
هي.

\textbf{2.} أبرد، بمقدار \qty{0.22}{\degreeCelsius}.

\textbf{3.} أدفأ، بمقدار \qty{0.98}{\degreeCelsius}.

\textbf{4.} $0.98 - (-0.22) = \qty{1.20}{\degreeCelsius}$.

\textbf{5.} $427 - 316 = \qty{111}{ppm}$.

\textbf{6.} $111 \div 316 \times 100 \approx 35~\%$.

\textbf{7.} $427 - 300 = \qty{127}{ppm}$.

\textbf{8.} يمتص ثاني أكسيد الكربون جزءًا من الأشعة تحت الحمراء التي يطلقها
السطح الدافئ ويعيد جزءًا منها نحو الأسفل. ومع وجود مزيد منه في الهواء، تعود
طاقة أكبر إلى السطح، فيسخن.

\textbf{9.} $2025 - 1993 = 32$ سنة؛ $32 \times 3.2 \approx
\qty{102}{mm}$، أي نحو \qty{10}{cm}.

\textbf{10.} ماء \omterm{def:g2:water-on-earth:ocean}{البحر} الذي يسخن يتمدد؛ والجليد على اليابسة (\omterm{def:g8:ice-on-earth:glacier}{الأنهار الجليدية}
\omterm{def:g8:ice-on-earth:ice-sheet}{والأغطية الجليدية}) يذوب في \omterm{def:g2:water-on-earth:ocean}{البحر}.

\textbf{11.} نعم، من الحجم نفسه: نحو \qty{1.2}{\degreeCelsius} مقابل
\qty{1.09}{\degreeCelsius}. ويختلفان قليلًا لأن الفترتين ليستا متطابقتين
(1880--1899 بدلًا من 1850--1900؛ و2015--2024 بدلًا من 2011--2020: والسنوات
الأحدث كانت أدفأ)، ولأن الهيئة الحكومية الدولية تجمع عدة سجلات، لا سجل ناسا
وحده.

\textbf{12.} «منذ نهاية القرن التاسع عشر، ارتفعت حرارة سطح الأرض بنحو
\textbf{\qty{1.2}{\degreeCelsius}}.»
\end{solution}
